\documentclass[11pt]{article}
\usepackage[margin=1in]{geometry}
\usepackage{amsmath}
\usepackage{booktabs}
\usepackage{graphicx}
\usepackage{hyperref}
\usepackage{enumitem}

\title{Benchmarking Simulated Annealing vs Hybrid QUBO Solvers in the Clearing Simulation}
\author{Project Technical Report}
\date{\today}

\begin{document}
\maketitle

\begin{abstract}
This report documents the single-run simulation mechanics and benchmarks two QUBO
solvers---Simulated Annealing (SA) and Hybrid---using 30 short simulations each.
We summarize model assumptions, daily calculations, and benchmarking constraints,
then analyze solver outcomes with caution due to the limited sample size and short
horizon. Results are derived from JSON logs in \texttt{benchmark\_results/} and
visualized using the corresponding PNG histograms.
\end{abstract}

\section{Introduction}
The clearing simulation models a multi-client, multi-asset clearinghouse where
daily market moves generate profits and losses, margins are computed via Expected
Shortfall, and a QUBO decision accepts or rejects proposed trades under a risk
budget. This report has two goals: (1) provide a scientific description of the
single-day algorithm and its assumptions; (2) compare SA and Hybrid solvers on a
fixed benchmark configuration.

\section{Simulation assumptions and configuration}
\subsection{Market and return generation}
Market dynamics use a three-regime Markov chain (calm, volatile, crisis), factor
returns, idiosyncratic noise, and optional jumps. These assumptions are encoded in
\texttt{clearing/market.py}:
\begin{itemize}[leftmargin=*]
  \item Transition matrix $P_z$ defines regime persistence.
  \item Returns are $r_t = B f_t + \epsilon_t$ with regime-dependent
  volatilities $\sigma_f[z_t]$ and $\sigma_e[z_t]$.
  \item Optional jumps add tail events with probability \texttt{jump\_p}.
\end{itemize}

\subsection{Scenario generation (RBM + quantization)}
Risk scenarios are sampled via a Restricted Boltzmann Machine (RBM) conditioned on
current regime $z_t$ and previous returns $r_{t-1}$ (see
\texttt{clearing/sampler.py}). Returns are quantized to $K$ bits per asset using
quantile bounds (see \texttt{clearing/quantization.py}). The sampler clamps
$(z_t, r_{t-1})$ and samples future returns $r_t$, producing scenario matrix
$R \in \mathbb{R}^{\Omega \times N}$.

\subsection{Portfolio, wealth, and collateral}
Each client holds a portfolio $P \in \mathbb{R}^{M \times N}$, wealth $W \in
\mathbb{R}^M$, and collateral $C \in \mathbb{R}^M$. Portfolios are initialized in
\texttt{scripts/run\_simulation.py} with a density mask and scaled random values.

\section{Daily simulation steps and calculations}
The daily loop is implemented in \texttt{clearing/simulation.py::simulate\_day}.
Phases correspond to the logger in \texttt{clearing/logger.py}.

\subsection{1. Market move}
Generate $z_t$ and $r_t$ (see \texttt{clearing/market.py::generate\_day}). PnL is
\[
\text{PnL}_i = \sum_{j=1}^N P_{ij} r_{t,j}
\]
Stress testing may apply an additional linear shock
$r_t \leftarrow r_t \cdot \texttt{stress\_return\_scale} + \texttt{stress\_return\_shift}$.

\subsection{2. Settlement and default loss}
Wealth and collateral are updated as:
\begin{align*}
W_{\text{settle}} &= W + \text{PnL} \\
F_{\text{pre}} &= (W - C) + \text{PnL} \\
C_{\text{after}} &= C + \min(F_{\text{pre}}, 0)
\end{align*}
Default loss accumulates if collateral becomes negative:
\[
\text{DL}_{\text{step}} = \max(-C_{\text{after}}, 0)
\]
See \texttt{clearing/simulation.py} for the settlement logic.

\subsection{3. Margining (Expected Shortfall)}
Scenario returns are sampled from the RBM and used to compute Expected Shortfall:
\begin{align*}
\text{PnL} &= R P^T \\
\text{Loss} &= -\text{PnL}
\end{align*}
ES is the mean of the worst $(1-\alpha)$ tail; VaR is the $\alpha$-quantile loss.
Implementation: \texttt{clearing/margin.py::margin\_es}.

\subsection{4. Default check and collateral top-up}
Clients must post collateral to meet the margin requirement. If available funds
are insufficient, the client defaults and is removed from trading:
\begin{align*}
F_{\text{avail}} &= \max(W_{\text{settle}} - C_{\text{after}}, 0) \\
\text{need} &= \max(M_{\text{req}} - C_{\text{after}}, 0)
\end{align*}
Defaulted clients are marked inactive and their portfolios are set to zero for
trading (\texttt{P\_liq}).

\subsection{5. Trade proposal}
Trades are generated using a momentum signal and exposure scaling in
\texttt{clearing/trading.py}:
\begin{align*}
\mu &= \texttt{base\_scale} \cdot \texttt{momentum\_weight} \cdot r_t \cdot (|P|+1)^{\gamma} \\
\sigma &= \texttt{noise\_scale} \cdot (\sqrt{|P|}+1) \\
\Delta P &= \mu + \sigma \cdot \epsilon,\quad \epsilon \sim \mathcal{N}(0,I)
\end{align*}
Zero positions are traded with probability \texttt{p\_trade\_when\_zero}. Trades are
clipped by \texttt{max\_abs\_trade}.

\subsection{6. QUBO clearing decision}
Tentative margins are computed for $P_{\text{tent}} = P_{\text{liq}} + \Delta P$,
with $\Delta M = M_{\text{req,tent}} - M_{\text{req,cur}}$. A Binary Quadratic Model
(BQM) selects clients to accept trades by minimizing a quadratic penalty that
enforces a risk budget $B$ and rewards risk-reducing trades. Implementation:
\texttt{clearing/qubo.py::build\_bqm\_accept\_clients}.

Solvers:
\begin{itemize}[leftmargin=*]
  \item \textbf{SA}: \texttt{dwave-neal} simulated annealing with
  \texttt{num\_reads} and \texttt{num\_sweeps} from \texttt{scripts/run\_simulation.py}.
  \item \textbf{Hybrid}: \texttt{dwave-system} LeapHybridSampler with a fixed
  time limit (set to 10 seconds in the benchmark harness) and per-call labels
  \texttt{ClSim\{run\}\_day\{d\}}.
\end{itemize}

\subsection{7. End of day}
Accepted trades update portfolios:
\[
P_{\text{next}} = P_{\text{liq}} + x \odot \Delta P
\]
Daily outputs include updated $P$, $W$, $C$, \texttt{alive}, default loss totals,
and QUBO statistics.

\section{Benchmark design}
\subsection{Inputs and settings}
The benchmark runs are recorded in:
\begin{itemize}[leftmargin=*]
  \item \texttt{benchmark\_results/benchmark\_sa\_20251230\_150952.json}
  \item \texttt{benchmark\_results/benchmark\_hybrid\_20251230\_163041.json}
\end{itemize}
Each solver is evaluated over 30 runs with $M=1000$ clients and $T=12$ days. The
number of assets is the default $N=500$ from \texttt{scripts/run\_simulation.py},
and the scenario count is the default $\Omega=1000$. Stress settings are fixed:
\texttt{stress}=True, \texttt{stress\_level}=1.0, \texttt{start\_state}=0,
\texttt{init\_position\_scale}=150, \texttt{init\_position\_density}=0.3,
\texttt{trade\_base\_scale}=10, \texttt{trade\_noise\_scale}=2,
\texttt{trade\_max\_abs}=100, \texttt{stress\_day}=3,
\texttt{stress\_return\_scale}=10, \texttt{stress\_return\_shift}=-0.05.
Seeds are incremented per run from base values (see JSON metadata).

\subsection{Metrics}
Each run records:
Total DL, Clients' gain (end minus start wealth), total defaulted clients,
mean QUBO size (variables + interactions), mean QUBO time, mean QUBO energy, and
total simulation time. Histogram overviews are in:
\begin{itemize}[leftmargin=*]
  \item \texttt{benchmark\_results/benchmark\_sa\_20251230\_150952.png}
  \item \texttt{benchmark\_results/benchmark\_hybrid\_20251230\_163041.png}
\end{itemize}

\subsection{Benchmarking limits}
Only 30 runs per solver are available, with a short 12-day horizon. This limits
statistical power and makes tail events noisy. Hybrid solver times include cloud
overhead and queue time; the requested time limit is 10 seconds per solve, but
observed solver time is higher. Results should be interpreted as indicative, not
conclusive.

\section{Results}
\subsection{Summary statistics}
Tables~\ref{tab:sa} and~\ref{tab:hybrid} summarize metrics (mean, median, min, max).

\begin{table}[h]
\centering
\caption{Simulated Annealing (30 runs, $M=1000$, $T=12$)}
\label{tab:sa}
\begin{tabular}{lrrrr}
\toprule
Metric & Mean & Median & Min & Max \\
\midrule
Total DL & 4332.45 & 1444.31 & 238.50 & 68991.84 \\
Clients' gain & 2710.15 & 352.15 & -23264.76 & 79903.50 \\
Total defaulted clients & 741.23 & 740.50 & 629.00 & 839.00 \\
Mean QUBO size & 47516.22 & 46588.92 & 35058.58 & 63258.33 \\
Mean QUBO time (s) & 4.39 & 4.26 & 3.52 & 5.46 \\
Mean QUBO energy & -332.16 & -337.16 & -455.30 & -206.13 \\
Total sim time (s) & 83.68 & 83.31 & 73.61 & 98.17 \\
\bottomrule
\end{tabular}
\end{table}

\begin{table}[h]
\centering
\caption{Hybrid Solver (30 runs, $M=1000$, $T=12$)}
\label{tab:hybrid}
\begin{tabular}{lrrrr}
\toprule
Metric & Mean & Median & Min & Max \\
\midrule
Total DL & 3608.13 & 1623.32 & 233.71 & 44984.77 \\
Clients' gain & 1851.29 & 2420.04 & -27501.65 & 99894.50 \\
Total defaulted clients & 742.37 & 744.50 & 645.00 & 852.00 \\
Mean QUBO size & 47980.41 & 47911.21 & 37466.92 & 60737.75 \\
Mean QUBO time (s) & 18.15 & 17.50 & 15.50 & 29.37 \\
Mean QUBO energy & -325.73 & -341.55 & -388.70 & -197.76 \\
Total sim time (s) & 247.82 & 238.68 & 217.08 & 385.60 \\
\bottomrule
\end{tabular}
\end{table}

\subsection{Histogram overview}
\begin{figure}[h]
\centering
\includegraphics[width=0.9\textwidth]{benchmark_results/benchmark_sa_20251230_150952.png}
\caption{SA histograms for all metrics (30 runs).}
\label{fig:sa_hist}
\end{figure}

\begin{figure}[h]
\centering
\includegraphics[width=0.9\textwidth]{benchmark_results/benchmark_hybrid_20251230_163041.png}
\caption{Hybrid histograms for all metrics (30 runs).}
\label{fig:hy_hist}
\end{figure}

\section{Discussion}
\subsection{Solver comparison}
Both solvers yield similar average defaulted-client counts and QUBO sizes, suggesting
the trade selection problem is of comparable scale across runs. Hybrid exhibits
substantially higher QUBO time (mean 18.15s vs 4.39s) and total simulation time
(mean 247.82s vs 83.68s), reflecting solver overhead and longer optimization.

For default loss (DL), Hybrid shows a lower mean but a higher median, while the
maximum DL is lower than SA (44984.77 vs 68991.84). This pattern may suggest
reduced extreme-tail events in Hybrid, but the small sample size and short horizon
make this inference weak. Clients' gain distributions are wide for both solvers,
with large positive and negative extremes, indicating strong sensitivity to random
initialization and market paths.

\subsection{High default-loss events}
DL distributions are highly skewed, with maxima orders of magnitude above medians.
This indicates occasional severe stress realizations, consistent with the stress
day at $t=3$ and amplified return scaling. Because only 30 runs are available,
these high-loss events dominate the mean and should not be interpreted as stable
tail statistics.

\section{Single-simulation limits}
Each run is only 12 days long with a fixed stress event, limiting insight into
long-horizon stability or rare-event frequency. Portfolio initialization is random
and sparse, and scenario generation depends on RBM sampling and quantizer settings.
Changing $N$, $K$, or the RBM/quantizer requires retraining to maintain consistency
between scenario sampling and return dynamics.

\section{Benchmarking limits}
The benchmark uses a single stress configuration and does not sweep budgets, trade
intensities, or market parameters. Hybrid solver performance is affected by external
latency and queue time. The results should be used for qualitative comparison
only and should not be over-interpreted as statistically robust.

\section{Conclusion}
In this benchmark, Hybrid is substantially slower but produces broadly similar
outcomes in default counts and QUBO size, with modest differences in default loss
statistics and energy. Given the limited sample size and short horizons, a more
extensive study (larger $T$, more runs, multiple stress profiles) is required
before drawing strong conclusions about solver quality or tail-risk mitigation.

\end{document}
